\section{Requirements Analysis} \label{sec:requirement_analysis}

This section outlines the requirements necessary to carry out the translation of the certified Imandra proof checker into the Rocq proof assistant. \autoref{sec:fun-requirements} presents the functional requirements of the system, while the non-functional requirements can be found in \autoref{sec:non-fun-requirements}.

\subsection{Functional Requirements} \label{sec:fun-requirements}

\paragraph*{FR1 --- Formalise the core proof system in Rocq.} The main pillar of the project is the formalisation of the proof certificates generated by Marabou. This includes a formal representation of the proof tree (FR1.1) and the DNN constraints (FR1.2).

\paragraph*{FR2 --- Parse proof certificates.} The system must be able to parse the proof certificates produced by Marabou and deserialise them into native Rocq data types (FR2.1). On the other hand, the system should handle and report any issues that arise during parsing (FR2.2).

\paragraph*{FR3 --- Verify proof certificates.} Once a certificate has been parsed, the system must determine whether it is valid by traversing the proof tree and checking each step (FR3.1).

\paragraph*{FR4 --- Support activation functions.} The checker should support the activation functions for which Marabou generates proof certificates. As a bare minimum, ReLU and the identity activation functions must be handled (FR4.1), as these are the only activation functions handled in the Imandra proof checker. Support for other piecewise-linear functions is desirable (FR4.2), while non-linear activation functions are considered out of scope for this project (FR4.3), since Marabou does not produce proofs for them.

\paragraph*{FR5 --- Formalise the DNN Farkas Lemma.} A certified proof of the DNN variant of the Farkas lemma must be developed in Rocq (FR5.1). This lemma is essential to establishing the soundness of the leaf-checking step.

\paragraph*{FR6 --- Formal proof of soundness.} The system must guarantee the soundness of its own implementation, that is, if the system returns UNSAT, then the corresponding DNN verification query genuinely has no solution (FR6.1).

\paragraph*{FR7 --- Integrate with external verification tools.} The system could be integrated with external tools such as Vehicle (FR7.1), which could make use of the certified proof checker as part of a broader verification pipeline.


\cref{tab:fun-requirements-sub} summarises the functional requirements
along with their MoSCoW priority classification.

\begin{table}[h]
\centering
\begin{tabularx}{\textwidth}{l X c }
\toprule
\textbf{ID} & \textbf{Description} & \textbf{MoSCoW} \\
\midrule
FR1.1 & Formal representation of a Marabou proof tree & Must \\
FR1.2 & Formal representation of DNN constraints & Must \\
FR2.1 & Parse Marabou certificates into native Rocq types & Must \\
FR2.2 & Detect and handle errors in proof certificates & Should \\
FR3.1 & Determine whether a proof certificate is valid & Must \\
FR4.1 & Support ReLU and identity activation functions & Must \\
FR4.2 & Support additional piecewise-linear activation functions & Could \\
FR4.3 & Support non-linear activation functions & Won't \\
FR5.1 & Certified proof of the DNN Farkas Lemma & Must \\
FR6.1 & Proof checker formal proof of soundness & Must \\
FR7.1 & Integration with external tools such as Vehicle & Could \\
\bottomrule
\end{tabularx}
\caption{Functional Requirements}
\label{tab:fun-requirements-sub}
\end{table}

\subsection{Non-Functional Requirements} \label{sec:non-fun-requirements}

\paragraph*{NFR1 --- Correctness.} All Rocq definitions and proofs must be accepted by the Rocq type checker (NFR1.1). In particular, no \texttt{admit} tactics should be used.

\paragraph*{NFR2 --- Reliability.} Verification results must be deterministic and reproducible across different environments (NFR2.1). Environment-specific configurations should be avoided at all costs.

\paragraph*{NFR3 --- Compatibility.} The implementation must be compatible with a stable release of Rocq and should avoid reliance on deprecated or unstable library features (NFR3.1).

\paragraph*{NFR4 --- Usability.} The codebase should be documented thoroughly to help the reader understand the codebase (NFR4.1), without the need for prior knowledge of the original Imandra implementation.

\cref{tab:non-fun-requirements-sub} summarises the non-functional requirements
along with their MoSCoW priority classification.

\begin{table}[h]
\centering
\begin{tabularx}{\textwidth}{l X c }
\toprule
\textbf{ID} & \textbf{Description} & \textbf{MoSCoW} \\
\midrule
NFR1.1 & All proofs accepted by Rocq type checker  & Must \\
NFR2.1 & Results are deterministic and reproducible & Must \\
NFR3.1 & Compatible with a stable Rocq release & Must \\
NFR4.1 & Codebase is documented for Rocq-familiar readers & Should \\
\bottomrule
\end{tabularx}
\caption{Non-Functional Requirements}
\label{tab:non-fun-requirements-sub}
\end{table}
